% Options for packages loaded elsewhere
\PassOptionsToPackage{unicode}{hyperref}
\PassOptionsToPackage{hyphens}{url}
\PassOptionsToPackage{dvipsnames,svgnames,x11names}{xcolor}
%
\documentclass[
  letterpaper,
  DIV=11,
  numbers=noendperiod]{scrartcl}

\usepackage{amsmath,amssymb}
\usepackage{iftex}
\ifPDFTeX
  \usepackage[T1]{fontenc}
  \usepackage[utf8]{inputenc}
  \usepackage{textcomp} % provide euro and other symbols
\else % if luatex or xetex
  \usepackage{unicode-math}
  \defaultfontfeatures{Scale=MatchLowercase}
  \defaultfontfeatures[\rmfamily]{Ligatures=TeX,Scale=1}
\fi
\usepackage{lmodern}
\ifPDFTeX\else  
    % xetex/luatex font selection
\fi
% Use upquote if available, for straight quotes in verbatim environments
\IfFileExists{upquote.sty}{\usepackage{upquote}}{}
\IfFileExists{microtype.sty}{% use microtype if available
  \usepackage[]{microtype}
  \UseMicrotypeSet[protrusion]{basicmath} % disable protrusion for tt fonts
}{}
\makeatletter
\@ifundefined{KOMAClassName}{% if non-KOMA class
  \IfFileExists{parskip.sty}{%
    \usepackage{parskip}
  }{% else
    \setlength{\parindent}{0pt}
    \setlength{\parskip}{6pt plus 2pt minus 1pt}}
}{% if KOMA class
  \KOMAoptions{parskip=half}}
\makeatother
\usepackage{xcolor}
\setlength{\emergencystretch}{3em} % prevent overfull lines
\setcounter{secnumdepth}{5}
% Make \paragraph and \subparagraph free-standing
\makeatletter
\ifx\paragraph\undefined\else
  \let\oldparagraph\paragraph
  \renewcommand{\paragraph}{
    \@ifstar
      \xxxParagraphStar
      \xxxParagraphNoStar
  }
  \newcommand{\xxxParagraphStar}[1]{\oldparagraph*{#1}\mbox{}}
  \newcommand{\xxxParagraphNoStar}[1]{\oldparagraph{#1}\mbox{}}
\fi
\ifx\subparagraph\undefined\else
  \let\oldsubparagraph\subparagraph
  \renewcommand{\subparagraph}{
    \@ifstar
      \xxxSubParagraphStar
      \xxxSubParagraphNoStar
  }
  \newcommand{\xxxSubParagraphStar}[1]{\oldsubparagraph*{#1}\mbox{}}
  \newcommand{\xxxSubParagraphNoStar}[1]{\oldsubparagraph{#1}\mbox{}}
\fi
\makeatother


\providecommand{\tightlist}{%
  \setlength{\itemsep}{0pt}\setlength{\parskip}{0pt}}\usepackage{longtable,booktabs,array}
\usepackage{calc} % for calculating minipage widths
% Correct order of tables after \paragraph or \subparagraph
\usepackage{etoolbox}
\makeatletter
\patchcmd\longtable{\par}{\if@noskipsec\mbox{}\fi\par}{}{}
\makeatother
% Allow footnotes in longtable head/foot
\IfFileExists{footnotehyper.sty}{\usepackage{footnotehyper}}{\usepackage{footnote}}
\makesavenoteenv{longtable}
\usepackage{graphicx}
\makeatletter
\def\maxwidth{\ifdim\Gin@nat@width>\linewidth\linewidth\else\Gin@nat@width\fi}
\def\maxheight{\ifdim\Gin@nat@height>\textheight\textheight\else\Gin@nat@height\fi}
\makeatother
% Scale images if necessary, so that they will not overflow the page
% margins by default, and it is still possible to overwrite the defaults
% using explicit options in \includegraphics[width, height, ...]{}
\setkeys{Gin}{width=\maxwidth,height=\maxheight,keepaspectratio}
% Set default figure placement to htbp
\makeatletter
\def\fps@figure{htbp}
\makeatother
% definitions for citeproc citations
\NewDocumentCommand\citeproctext{}{}
\NewDocumentCommand\citeproc{mm}{%
  \begingroup\def\citeproctext{#2}\cite{#1}\endgroup}
\makeatletter
 % allow citations to break across lines
 \let\@cite@ofmt\@firstofone
 % avoid brackets around text for \cite:
 \def\@biblabel#1{}
 \def\@cite#1#2{{#1\if@tempswa , #2\fi}}
\makeatother
\newlength{\cslhangindent}
\setlength{\cslhangindent}{1.5em}
\newlength{\csllabelwidth}
\setlength{\csllabelwidth}{3em}
\newenvironment{CSLReferences}[2] % #1 hanging-indent, #2 entry-spacing
 {\begin{list}{}{%
  \setlength{\itemindent}{0pt}
  \setlength{\leftmargin}{0pt}
  \setlength{\parsep}{0pt}
  % turn on hanging indent if param 1 is 1
  \ifodd #1
   \setlength{\leftmargin}{\cslhangindent}
   \setlength{\itemindent}{-1\cslhangindent}
  \fi
  % set entry spacing
  \setlength{\itemsep}{#2\baselineskip}}}
 {\end{list}}
\usepackage{calc}
\newcommand{\CSLBlock}[1]{\hfill\break\parbox[t]{\linewidth}{\strut\ignorespaces#1\strut}}
\newcommand{\CSLLeftMargin}[1]{\parbox[t]{\csllabelwidth}{\strut#1\strut}}
\newcommand{\CSLRightInline}[1]{\parbox[t]{\linewidth - \csllabelwidth}{\strut#1\strut}}
\newcommand{\CSLIndent}[1]{\hspace{\cslhangindent}#1}

\DeclareMathAlphabet{\mathcal}{OMS}{cmsy}{m}{n}
\KOMAoption{captions}{tableheading}
\makeatletter
\@ifpackageloaded{caption}{}{\usepackage{caption}}
\AtBeginDocument{%
\ifdefined\contentsname
  \renewcommand*\contentsname{Table of contents}
\else
  \newcommand\contentsname{Table of contents}
\fi
\ifdefined\listfigurename
  \renewcommand*\listfigurename{List of Figures}
\else
  \newcommand\listfigurename{List of Figures}
\fi
\ifdefined\listtablename
  \renewcommand*\listtablename{List of Tables}
\else
  \newcommand\listtablename{List of Tables}
\fi
\ifdefined\figurename
  \renewcommand*\figurename{Figure}
\else
  \newcommand\figurename{Figure}
\fi
\ifdefined\tablename
  \renewcommand*\tablename{Table}
\else
  \newcommand\tablename{Table}
\fi
}
\@ifpackageloaded{float}{}{\usepackage{float}}
\floatstyle{ruled}
\@ifundefined{c@chapter}{\newfloat{codelisting}{h}{lop}}{\newfloat{codelisting}{h}{lop}[chapter]}
\floatname{codelisting}{Listing}
\newcommand*\listoflistings{\listof{codelisting}{List of Listings}}
\makeatother
\makeatletter
\makeatother
\makeatletter
\@ifpackageloaded{caption}{}{\usepackage{caption}}
\@ifpackageloaded{subcaption}{}{\usepackage{subcaption}}
\makeatother

\ifLuaTeX
  \usepackage{selnolig}  % disable illegal ligatures
\fi
\usepackage{bookmark}

\IfFileExists{xurl.sty}{\usepackage{xurl}}{} % add URL line breaks if available
\urlstyle{same} % disable monospaced font for URLs
\hypersetup{
  pdftitle={OpenREML Technical Documentation},
  pdfauthor={Patrick Li},
  colorlinks=true,
  linkcolor={blue},
  filecolor={Maroon},
  citecolor={Blue},
  urlcolor={Blue},
  pdfcreator={LaTeX via pandoc}}


\title{OpenREML Technical Documentation}
\author{Patrick Li}
\date{}

\begin{document}
\maketitle

\renewcommand*\contentsname{Table of contents}
{
\hypersetup{linkcolor=}
\setcounter{tocdepth}{3}
\tableofcontents
}

\section{Gaussian Linear Mixed Model
Formulation}\label{gaussian-linear-mixed-model-formulation}

Consider the linear mixed model \begin{equation}
\symbf{Y} = \symbf{X}\symbf{\beta} + \symbf{Z}\symbf{b} + \symbf{\varepsilon},
\end{equation} where:

\begin{itemize}
\tightlist
\item
  \(\symbf{Y} \in \mathbb{R}^{n \times 1}\) is the response vector,
\item
  \(\symbf{X} \in \mathbb{R}^{n \times p}\) is the fixed-effects design
  matrix,
\item
  \(\symbf{\beta} \in \mathbb{R}^{p \times 1}\) is the vector of
  fixed-effect coefficients,
\item
  \(\symbf{Z} \in \mathbb{R}^{n \times q}\) is the random-effects design
  matrix,
\item
  \(\symbf{b} \in \mathbb{R}^{q \times 1}\) is the vector of random
  effects,
\item
  \(\symbf{\varepsilon} \in \mathbb{R}^{n \times 1}\) is the vector of
  residual errors.
\end{itemize}

We assume \begin{equation}
\symbf{b} \sim \mathcal{N}(\symbf{0}, \symbf{G}), \quad
\symbf{\varepsilon} \sim \mathcal{N}(\symbf{0}, \symbf{R}), \quad
\symbf{b} \perp \symbf{\varepsilon},
\end{equation} where \(\symbf{G} \in \mathbb{R}^{q \times q}\) and
\(\symbf{R} \in \mathbb{R}^{n \times n}\) are covariance matrices.

Let \(\symbf{\theta}\) denote the collection of variance components that
parameterise \(\symbf{G}\) and \(\symbf{R}\), i.e., \begin{equation}
\symbf{G} = \symbf{G}(\symbf{\theta}), \quad
\symbf{R} = \symbf{R}(\symbf{\theta}).
\end{equation}

Under these assumptions, the marginal distribution of \(\symbf{Y}\) is
given by \begin{equation}
\symbf{Y} \sim \mathcal{N}\left(\symbf{X}\symbf{\beta},\; \symbf{\Sigma}(\symbf{\theta})\right),
\end{equation} where \begin{equation}
\symbf{\Sigma}(\symbf{\theta}) = \symbf{Z}\symbf{G}(\symbf{\theta})\symbf{Z}^\top + \symbf{R}(\symbf{\theta}).
\end{equation}

\section{BLUE}\label{blue}

Given \(\symbf{\theta}\), the Best Linear Unbiased Estimator (BLUE) of
\(\symbf{\beta}\) is the linear unbiased estimator that minimises the
variance among all such estimators. Under the linear mixed model with
covariance matrix \(\symbf{\Sigma}(\symbf{\theta})\), the BLUE of
\(\symbf{\beta}\) is given by \begin{equation}
\widehat{\symbf{\beta}}(\symbf{\theta})
=
\left(\symbf{X}^\top \symbf{\Sigma}(\symbf{\theta})^{-1} \symbf{X}\right)^{-1}
\symbf{X}^\top \symbf{\Sigma}(\symbf{\theta})^{-1} \symbf{Y}.
\end{equation}

This estimator is the generalised least squares (GLS) estimator of
\(\symbf{\beta}\).

\section{BLUP}\label{blup}

Given \(\symbf{\theta}\), we often wish to predict the random effects
\(\symbf{b}\). The Best Linear Unbiased Predictor (BLUP) of
\(\symbf{b}\) is given by:

\begin{equation}
\widehat{\symbf{b}}(\symbf{\theta}) = \symbf{G}(\symbf{\theta})\symbf{Z}^\top \symbf{\Sigma}^{-1}(\symbf{\theta}) (\symbf{Y} - \symbf{X}\widehat{\symbf{\beta}}).
\label{eq:blup}
\end{equation}

\section{ML Estimation}\label{ml-estimation}

Maximum likelihood (ML) estimation jointly estimates the fixed effects
\(\symbf{\beta}\) and variance components \(\symbf{\theta}\) by
maximising the marginal likelihood of the observed data. After
integrating out the random effects \(\symbf{b}\) from the joint model,
the marginal distribution of \(\symbf{Y}\) is \begin{equation}
\symbf{Y} \sim \mathcal{N}\!\left(\symbf{X}\symbf{\beta},\; \symbf{\Sigma}(\symbf{\theta})\right),
\end{equation} where
\(\symbf{\Sigma}(\symbf{\theta}) = \symbf{Z}\symbf{G}(\symbf{\theta})\symbf{Z}^\top + \symbf{R}(\symbf{\theta})\).
The corresponding log-likelihood is given by \begin{equation}
\ell(\symbf{\beta}, \symbf{\theta})
=
-\frac{1}{2}\left[
n\log(2\pi)
+ \log\left|\symbf{\Sigma}(\symbf{\theta})\right|
+ (\symbf{Y} - \symbf{X}\symbf{\beta})^\top
\symbf{\Sigma}(\symbf{\theta})^{-1}
(\symbf{Y} - \symbf{X}\symbf{\beta})
\right].
\end{equation}

For fixed \(\symbf{\theta}\), the log-likelihood is quadratic in
\(\symbf{\beta}\), and the maximiser can be obtained in closed form by
setting \(\partial \ell / \partial \symbf{\beta} = \symbf{0}\). This
yields the generalised least squares (GLS) estimator \begin{equation}
\widehat{\symbf{\beta}}(\symbf{\theta})
=
\left(
\symbf{X}^\top \symbf{\Sigma}(\symbf{\theta})^{-1} \symbf{X}
\right)^{-1}
\symbf{X}^\top \symbf{\Sigma}(\symbf{\theta})^{-1} \symbf{Y}.
\end{equation} Substituting this expression back into the log-likelihood
eliminates \(\symbf{\beta}\) and produces the profile log-likelihood
\begin{equation}
\ell_p(\symbf{\theta})
=
-\frac{1}{2}\left[
n\log(2\pi)
+ \log\left|\symbf{\Sigma}(\symbf{\theta})\right|
+ \symbf{Y}^\top \symbf{P}(\symbf{\theta})\, \symbf{Y}
\right],
\end{equation} where \begin{equation}
\symbf{P}(\symbf{\theta})
=
\symbf{\Sigma}(\symbf{\theta})^{-1}
-
\symbf{\Sigma}(\symbf{\theta})^{-1}\symbf{X}
\left(
\symbf{X}^\top\symbf{\Sigma}(\symbf{\theta})^{-1}\symbf{X}
\right)^{-1}
\symbf{X}^\top\symbf{\Sigma}(\symbf{\theta})^{-1}
\end{equation} is the residual projection matrix. The resulting
optimisation problem depends only on \(\symbf{\theta}\).

Since \(\ell_p(\symbf{\theta})\) is nonlinear in \(\symbf{\theta}\) and
does not admit a closed-form maximiser, numerical optimisation methods
are required. These methods rely on the score vector and a curvature
approximation. Differentiating the profile log-likelihood with respect
to a variance component \(\theta_k\) gives \begin{equation}
\frac{\partial \ell_p}{\partial \theta_k}
=
-\frac{1}{2}\left[
\operatorname{tr}\!\left(
\symbf{\Sigma}^{-1}\frac{\partial \symbf{\Sigma}}{\partial \theta_k}
\right)
-
\symbf{Y}^\top \symbf{P}
\frac{\partial \symbf{\Sigma}}{\partial \theta_k}
\symbf{P}\,\symbf{Y}
\right],
\end{equation} where the identity
\(\partial \log|\symbf{\Sigma}|/\partial \theta_k
=
\operatorname{tr}(\symbf{\Sigma}^{-1}\partial \symbf{\Sigma}/\partial \theta_k)\)
has been used, together with \(\partial \symbf{P}/\partial \theta_k
=
-\symbf{P}(\partial \symbf{\Sigma}/\partial \theta_k)\symbf{P}\).

A common curvature approximation is given by the Fisher information
matrix, whose \((k,j)\) entry is \begin{equation}
\mathcal{I}_{kj}(\symbf{\theta})
=
\frac{1}{2}\operatorname{tr}\!\left(
\symbf{\Sigma}^{-1}\frac{\partial\symbf{\Sigma}}{\partial\theta_k}
\symbf{\Sigma}^{-1}\frac{\partial\symbf{\Sigma}}{\partial\theta_j}
\right).
\end{equation} Starting from an initial value \(\symbf{\theta}^{(0)}\),
the Fisher scoring iteration updates \(\symbf{\theta}\) according to
\begin{equation}
\symbf{\theta}^{(t+1)}
=
\symbf{\theta}^{(t)}
+
\mathcal{I}\!\left(\symbf{\theta}^{(t)}\right)^{-1}
\symbf{s}\!\left(\symbf{\theta}^{(t)}\right),
\end{equation} where \(\symbf{s}(\symbf{\theta})\) denotes the score
vector. In practice, variance components are typically reparameterised
(e.g.~on a logarithmic scale) to ensure positivity during optimisation.
Iterations are continued until convergence, for example when
\(\|\symbf{\theta}^{(t+1)} - \symbf{\theta}^{(t)}\| < \varepsilon\) for
a prescribed tolerance \(\varepsilon > 0\).

Upon convergence, the ML estimate of \(\symbf{\theta}\) is
\(\widehat{\symbf{\theta}}_{ML}\), and the ML estimate of
\(\symbf{\beta}_{ML}\) is recovered analytically as \begin{equation}
    \widehat{\symbf{\beta}}_{ML}
    = \widehat{\symbf{\beta}}\!\left(\widehat{\symbf{\theta}}_{ML}\right)
    =
    \left(
        \symbf{X}^\top
        \symbf{\Sigma}(\widehat{\symbf{\theta}}_{ML})^{-1}
        \symbf{X}
    \right)^{-1}
    \symbf{X}^\top
    \symbf{\Sigma}(\widehat{\symbf{\theta}}_{ML})^{-1}
    \symbf{Y}.
\end{equation}

\noindent\textbf{Remark.} ML treats \(\symbf{\beta}\) and
\(\symbf{\theta}\) symmetrically in the likelihood and is therefore
consistent as \(n \to \infty\). However, in finite samples ML
underestimates variance components because it does not account for the
degrees of freedom consumed by estimating \(\symbf{\beta}\). This bias
motivates restricted maximum likelihood (REML), which maximises a
likelihood formed from residual contrasts that are orthogonal to the
column space of \(\symbf{X}\), thereby adjusting for the estimation of
fixed effects when recovering \(\symbf{\theta}\).

\section{REML Estimation}\label{reml-estimation}

Restricted maximum likelihood (REML) constructs a likelihood based on
linear combinations of \(\symbf{Y}\) that are free of \(\symbf{\beta}\).
Let \(\symbf{A}\) be an \((n-p)\times n\) matrix satisfying
\(\symbf{A}\symbf{X}=\symbf{0}\) and \(rank(\symbf{A})=n-p\). Then
\begin{equation}
\symbf{AY} \sim \mathcal{N}\big(\symbf{0},\; \symbf{A}\symbf{\Sigma}(\symbf{\theta})\symbf{A}^\top\big),
\end{equation} and the corresponding log-likelihood is \begin{equation}
\ell_R(\symbf{\theta})
=
-\frac{1}{2}\left[
(n-p)\log(2\pi)
+ \log\big|\symbf{A}\symbf{\Sigma}(\symbf{\theta})\symbf{A}^\top\big|
+ (\symbf{AY})^\top \big(\symbf{A}\symbf{\Sigma}(\symbf{\theta})\symbf{A}^\top\big)^{-1} (\symbf{AY})
\right].
\end{equation}

Using standard matrix identities, this expression can be rewritten as
\begin{equation}
\boxed{
\ell_R(\symbf{\theta})
=
-\frac{1}{2}\left[
\log\big|\symbf{\Sigma}(\symbf{\theta})\big|
+ \log\big|\symbf{X}^\top \symbf{\Sigma}(\symbf{\theta})^{-1} \symbf{X}\big|
+ \symbf{Y}^\top \symbf{P}(\symbf{\theta}) \symbf{Y}
\right]
+ C
},
\end{equation} where \begin{equation}
\boxed{
\symbf{P}(\symbf{\theta})
=
\symbf{\Sigma}(\symbf{\theta})^{-1}
-
\symbf{\Sigma}(\symbf{\theta})^{-1}\symbf{X}
\left(
\symbf{X}^\top \symbf{\Sigma}(\symbf{\theta})^{-1}\symbf{X}
\right)^{-1}
\symbf{X}^\top \symbf{\Sigma}(\symbf{\theta})^{-1}
},
\end{equation} and \(C\) is a constant independent of
\(\symbf{\theta}\). Since additive constants do not affect maximisation
with respect to \(\symbf{\theta}\), the REML log-likelihood is taken to
be the above expression up to an additive constant.

\subsection{Score Function}\label{score-function}

Let \(\theta_k\) denote the \(k\)-th component of \(\symbf{\theta}\),
and define \begin{equation}
\dot{\symbf{\Sigma}}_k
=
\frac{\partial \symbf{\Sigma}(\symbf{\theta})}{\partial \theta_k}.
\end{equation} Starting from \begin{equation}
\ell_R(\symbf{\theta})
=
-\frac{1}{2}\left[
\log|\symbf{\Sigma}|
+ \log\big|\symbf{X}^\top \symbf{\Sigma}^{-1}\symbf{X}\big|
+ \symbf{Y}^\top \symbf{P}\symbf{Y}
\right]
+ C,
\end{equation} we differentiate each term with respect to \(\theta_k\).

\subsubsection*{Derivative of $\log|\symbf{\Sigma}|$}

\begin{equation}
\frac{\partial}{\partial \theta_k}\log|\symbf{\Sigma}|
=
\operatorname{tr}\!\big(\symbf{\Sigma}^{-1}\dot{\symbf{\Sigma}}_k\big).
\end{equation}

\subsubsection*{Derivative of $\log|\symbf{X}^\top \symbf{\Sigma}^{-1}\symbf{X}|$}

Set \(\symbf{M} = \symbf{X}^\top \symbf{\Sigma}^{-1}\symbf{X}\). Then
\begin{equation}
\frac{\partial}{\partial \theta_k}\log|\symbf{M}|
=
\operatorname{tr}\!\big(\symbf{M}^{-1} \dot{\symbf{M}}_k\big),
\end{equation} where \begin{equation}
\dot{\symbf{M}}_k
=
\frac{\partial}{\partial \theta_k}
\big(\symbf{X}^\top \symbf{\Sigma}^{-1}\symbf{X}\big)
=
-\,\symbf{X}^\top \symbf{\Sigma}^{-1}
\dot{\symbf{\Sigma}}_k
\symbf{\Sigma}^{-1}\symbf{X}.
\end{equation} Hence, \begin{equation}
\frac{\partial}{\partial \theta_k}
\log\big|\symbf{X}^\top \symbf{\Sigma}^{-1}\symbf{X}\big|
=
-\,\operatorname{tr}\!\Big(
\big(\symbf{X}^\top \symbf{\Sigma}^{-1}\symbf{X}\big)^{-1}
\symbf{X}^\top \symbf{\Sigma}^{-1}
\dot{\symbf{\Sigma}}_k
\symbf{\Sigma}^{-1}\symbf{X}
\Big).
\end{equation}

\subsubsection{\texorpdfstring{Derivative of
\(\symbf{Y}^\top \symbf{P}\symbf{Y}\)}{Derivative of \textbackslash symbf\{Y\}\^{}\textbackslash top \textbackslash symbf\{P\}\textbackslash symbf\{Y\}}}\label{derivative-of-symbfytop-symbfpsymbfy}

\begin{equation}
\frac{\partial}{\partial \theta_k}
\big(\symbf{Y}^\top \symbf{P}\symbf{Y}\big)
=
\symbf{Y}^\top
\frac{\partial \symbf{P}}{\partial \theta_k}
\symbf{Y}.
\end{equation} Using \begin{equation}
\frac{\partial \symbf{P}}{\partial \theta_k}
=
-\,\symbf{P}\,\dot{\symbf{\Sigma}}_k\,\symbf{P},
\end{equation} we obtain \begin{equation}
\frac{\partial}{\partial \theta_k}
\big(\symbf{Y}^\top \symbf{P}\symbf{Y}\big)
=
-\,\symbf{Y}^\top \symbf{P}\,\dot{\symbf{\Sigma}}_k\,\symbf{P}\,\symbf{Y}.
\end{equation}

\subsubsection{Compact Form}\label{compact-form}

Combining the three components and noting that \(C\) does not depend on
\(\symbf{\theta}\), the score function is \begin{equation}
S_k(\symbf{\theta})
=
\frac{\partial \ell_R(\symbf{\theta})}{\partial \theta_k}
=
-\frac{1}{2}
\left[
\operatorname{tr}\!\big(\symbf{\Sigma}^{-1}\dot{\symbf{\Sigma}}_k\big)
-
\operatorname{tr}\!\Big(
\big(\symbf{X}^\top \symbf{\Sigma}^{-1}\symbf{X}\big)^{-1}
\symbf{X}^\top \symbf{\Sigma}^{-1}
\dot{\symbf{\Sigma}}_k
\symbf{\Sigma}^{-1}\symbf{X}
\Big)
-
\symbf{Y}^\top \symbf{P}\,\dot{\symbf{\Sigma}}_k\,\symbf{P}\,\symbf{Y}
\right].
\end{equation}

To obtain a more compact form, recall the identity \begin{equation}
\symbf{P}
=
\symbf{\Sigma}^{-1}
-
\symbf{\Sigma}^{-1}\symbf{X}
\left(
\symbf{X}^\top \symbf{\Sigma}^{-1}\symbf{X}
\right)^{-1}
\symbf{X}^\top \symbf{\Sigma}^{-1}.
\end{equation} Hence, \begin{equation}
\operatorname{tr}\!\big(\symbf{P}\,\dot{\symbf{\Sigma}}_k\big)
=
\operatorname{tr}\!\big(\symbf{\Sigma}^{-1}\dot{\symbf{\Sigma}}_k\big)
-
\operatorname{tr}\!\Big(
\symbf{\Sigma}^{-1}\symbf{X}
\left(
\symbf{X}^\top \symbf{\Sigma}^{-1}\symbf{X}
\right)^{-1}
\symbf{X}^\top \symbf{\Sigma}^{-1}
\dot{\symbf{\Sigma}}_k
\Big).
\end{equation} Using the cyclic property of the trace,
\(\operatorname{tr}(\symbf{D}\symbf{A}\symbf{B}\symbf{C})
=
\operatorname{tr}(\symbf{A}\symbf{B}\symbf{C}\symbf{D})\), the second
term becomes \begin{equation}
\operatorname{tr}\!\Big(
\symbf{\Sigma}^{-1}\symbf{X}
\left(
\symbf{X}^\top \symbf{\Sigma}^{-1}\symbf{X}
\right)^{-1}
\symbf{X}^\top \symbf{\Sigma}^{-1}
\dot{\symbf{\Sigma}}_k
\Big)
=
\operatorname{tr}\!\Big(
\left(
\symbf{X}^\top \symbf{\Sigma}^{-1}\symbf{X}
\right)^{-1}
\symbf{X}^\top \symbf{\Sigma}^{-1}
\dot{\symbf{\Sigma}}_k
\symbf{\Sigma}^{-1}\symbf{X}
\Big).
\end{equation} Therefore, \begin{equation}
\operatorname{tr}\!\big(\symbf{P}\,\dot{\symbf{\Sigma}}_k\big)
=
\operatorname{tr}\!\big(\symbf{\Sigma}^{-1}\dot{\symbf{\Sigma}}_k\big)
-
\operatorname{tr}\!\Big(
\big(\symbf{X}^\top \symbf{\Sigma}^{-1}\symbf{X}\big)^{-1}
\symbf{X}^\top \symbf{\Sigma}^{-1}
\dot{\symbf{\Sigma}}_k
\symbf{\Sigma}^{-1}\symbf{X}
\Big).
\end{equation} Substituting this identity into the previous expression
for \(S_k(\symbf{\theta})\) yields \begin{equation}
\label{eq-score}
\boxed{
S_k(\symbf{\theta})
=
\frac{1}{2}
\left[
\symbf{Y}^\top \symbf{P}\,\dot{\symbf{\Sigma}}_k\,\symbf{P}\,\symbf{Y}
-
\operatorname{tr}\!\big(\symbf{P}\,\dot{\symbf{\Sigma}}_k\big)
\right]
}.
\end{equation}

\subsection{Average Information
Matrix}\label{average-information-matrix}

The score function for \(\theta_k\) is \begin{equation}
S_k(\symbf{\theta})
=
\frac{1}{2}
\left[
\symbf{Y}^\top \symbf{P}\,\dot{\symbf{\Sigma}}_k\,\symbf{P}\,\symbf{Y}
-
\operatorname{tr}\!\big(\symbf{P}\,\dot{\symbf{\Sigma}}_k\big)
\right],
\end{equation} where \begin{equation}
\dot{\symbf{\Sigma}}_k
=
\frac{\partial \symbf{\Sigma}(\symbf{\theta})}{\partial \theta_k}.
\end{equation}

\subsubsection{Observed Information}\label{observed-information}

Differentiating \(S_k\) with respect to \(\theta_j\) gives
\begin{equation}
-\frac{\partial^2 \ell_R}{\partial \theta_k \partial \theta_j}
=
-\frac{1}{2}
\left[
\symbf{Y}^\top
\frac{\partial}{\partial \theta_j}
\big(\symbf{P}\,\dot{\symbf{\Sigma}}_k\,\symbf{P}\big)
\symbf{Y}
-
\frac{\partial}{\partial \theta_j}
\operatorname{tr}\!\big(\symbf{P}\,\dot{\symbf{\Sigma}}_k\big)
\right].
\end{equation}

For the quadratic term, using the product rule and
\(\partial \symbf{P}/\partial \theta_j = -\symbf{P}\,\dot{\symbf{\Sigma}}_j\,\symbf{P}\),
we obtain \begin{equation}
\frac{\partial}{\partial \theta_j}
\big(\symbf{P}\,\dot{\symbf{\Sigma}}_k\,\symbf{P}\big)
=
-\,\symbf{P}\,\dot{\symbf{\Sigma}}_j\,\symbf{P}\,\dot{\symbf{\Sigma}}_k\,\symbf{P}
-
\symbf{P}\,\dot{\symbf{\Sigma}}_k\,\symbf{P}\,\dot{\symbf{\Sigma}}_j\,\symbf{P}
+
\symbf{P}\,\ddot{\symbf{\Sigma}}_{kj}\,\symbf{P}.
\end{equation}

For the trace term, \begin{equation}
\frac{\partial}{\partial \theta_j}
\operatorname{tr}\!\big(\symbf{P}\,\dot{\symbf{\Sigma}}_k\big)
=
-\,\operatorname{tr}\!\big(
\symbf{P}\,\dot{\symbf{\Sigma}}_j\,\symbf{P}\,\dot{\symbf{\Sigma}}_k
\big)
+
\operatorname{tr}\!\big(
\symbf{P}\,\ddot{\symbf{\Sigma}}_{kj}
\big).
\end{equation}

Putting everything together yields the observed information:
\begin{equation}
\mathcal{J}_{kj}(\symbf{\theta}) = -\frac{\partial^2 \ell_R}{\partial \theta_k \partial \theta_j}
=
\symbf{Y}^\top
\symbf{P}\,\dot{\symbf{\Sigma}}_k\,\symbf{P}\,\dot{\symbf{\Sigma}}_j\,\symbf{P}
\symbf{Y}
-\frac{1}{2}
\operatorname{tr}\!\big(
\symbf{P}\,\dot{\symbf{\Sigma}}_j\,\symbf{P}\,\dot{\symbf{\Sigma}}_k
\big)
+\frac{1}{2}
\operatorname{tr}\!\big(
\symbf{P}\,\ddot{\symbf{\Sigma}}_{kj}
\big)
- \frac{1}{2} \symbf{Y}^\top \symbf{P}\,\ddot{\symbf{\Sigma}}_{kj}\,\symbf{P}\symbf{Y}.
\end{equation}

\subsubsection{Expected Information}\label{expected-information}

We now derive the expected Fisher information matrix by taking
expectations of the negative second derivatives. Recall that under the
model, \(\mathbb{E}[\symbf{Y}] = \symbf{X}\symbf{\beta}\). However, a
fundamental property of the projection matrix \(\symbf{P}\) is that
\(\symbf{P}\symbf{X} = \symbf{0}\), which implies
\(\symbf{P}\symbf{Y} = \symbf{P}(\symbf{Y} - \symbf{X}\symbf{\beta})\).
Consequently, when computing expectations of expressions involving
\(\symbf{P}\), we may without loss of generality treat \(\symbf{Y}\) as
centered at zero. We therefore take
\(\mathbb{E}[\symbf{Y}] = \symbf{0}\) and
\(\mathbb{E}[\symbf{Y}\symbf{Y}^\top] = \symbf{\Sigma}\) for the
remainder of this derivation.

The observed information matrix is given by \begin{equation}
-\frac{\partial^2 \ell_R}{\partial \theta_k \partial \theta_j}
=
\underbrace{\symbf{Y}^\top \symbf{P}\,\dot{\symbf{\Sigma}}_k\,\symbf{P}\,\dot{\symbf{\Sigma}}_j\,\symbf{P} \symbf{Y}}_{T_1}
-\frac12 \underbrace{\operatorname{tr}\!\big( \symbf{P}\,\dot{\symbf{\Sigma}}_j\,\symbf{P}\,\dot{\symbf{\Sigma}}_k \big)}_{T_2}
+\frac12 \underbrace{\operatorname{tr}\!\big( \symbf{P}\,\ddot{\symbf{\Sigma}}_{kj} \big)}_{T_3}
- \frac12 \underbrace{\symbf{Y}^\top \symbf{P}\,\ddot{\symbf{\Sigma}}_{kj}\,\symbf{P}\symbf{Y}}_{T_4}.
\end{equation} We evaluate the expectation of each term in turn.

\paragraph{\texorpdfstring{Expectation of
\(T_1\)}{Expectation of T\_1}}\label{expectation-of-t_1}

The quantity \(T_1\) is a quadratic form in \(\symbf{Y}\). For a
zero-mean random vector with covariance \(\symbf{\Sigma}\), the identity
\(\mathbb{E}[\symbf{Y}^\top \symbf{A} \symbf{Y}] = \operatorname{tr}(\symbf{A} \symbf{\Sigma})\)
holds for any conformable matrix \(\symbf{A}\). Applying this with
\(\symbf{A} = \symbf{P}\,\dot{\symbf{\Sigma}}_k\,\symbf{P}\,\dot{\symbf{\Sigma}}_j\,\symbf{P}\)
yields \begin{equation}
\mathbb{E}[T_1] = \operatorname{tr}\!\big( \symbf{P}\,\dot{\symbf{\Sigma}}_k\,\symbf{P}\,\dot{\symbf{\Sigma}}_j\,\symbf{P} \, \symbf{\Sigma} \big).
\end{equation} A key identity satisfied by \(\symbf{P}\) is
\(\symbf{P}\symbf{\Sigma}\symbf{P} = \symbf{P}\), which can be verified
by direct substitution of the definition
\(\symbf{P} = \symbf{\Sigma}^{-1} - \symbf{\Sigma}^{-1}\symbf{X}(\symbf{X}^\top\symbf{\Sigma}^{-1}\symbf{X})^{-1}\symbf{X}^\top\symbf{\Sigma}^{-1}\).
Using the cyclic property of the trace, \begin{align}
\mathbb{E}[T_1] &= \operatorname{tr}\!\big( \symbf{P}\,\dot{\symbf{\Sigma}}_k\,\symbf{P}\,\dot{\symbf{\Sigma}}_j\, (\symbf{P}\symbf{\Sigma}) \big) \\
&= \operatorname{tr}\!\big( \dot{\symbf{\Sigma}}_k\,\symbf{P}\,\dot{\symbf{\Sigma}}_j\,\symbf{P}\symbf{\Sigma}\symbf{P} \big) \quad \text{(cyclic permutation)} \\
&= \operatorname{tr}\!\big( \dot{\symbf{\Sigma}}_k\,\symbf{P}\,\dot{\symbf{\Sigma}}_j\,\symbf{P} \big) \quad \text{(since $\symbf{P}\symbf{\Sigma}\symbf{P} = \symbf{P}$)} \\
&= \operatorname{tr}\!\big( \symbf{P}\,\dot{\symbf{\Sigma}}_k\,\symbf{P}\,\dot{\symbf{\Sigma}}_j \big) \quad \text{(cyclic permutation again)}.
\end{align} Thus \begin{equation}
\mathbb{E}[T_1] = \operatorname{tr}\!\big( \symbf{P}\,\dot{\symbf{\Sigma}}_k\,\symbf{P}\,\dot{\symbf{\Sigma}}_j \big).
\end{equation}

\paragraph{\texorpdfstring{Expectation of
\(T_2\)}{Expectation of T\_2}}\label{expectation-of-t_2}

The term \(T_2\) contains no random quantities; it is a deterministic
function of the parameters. Hence \begin{equation}
\mathbb{E}[T_2] = \operatorname{tr}\!\big( \symbf{P}\,\dot{\symbf{\Sigma}}_j\,\symbf{P}\,\dot{\symbf{\Sigma}}_k \big).
\end{equation}

\paragraph{\texorpdfstring{Expectation of
\(T_3\)}{Expectation of T\_3}}\label{expectation-of-t_3}

Similarly, \(T_3\) is non-random, giving \begin{equation}
\mathbb{E}[T_3] = \operatorname{tr}\!\big( \symbf{P}\,\ddot{\symbf{\Sigma}}_{kj} \big).
\end{equation}

\paragraph{\texorpdfstring{Expectation of
\(T_4\)}{Expectation of T\_4}}\label{expectation-of-t_4}

The term \(T_4\) is again a quadratic form. Applying the same identity
as for \(T_1\), \begin{equation}
\mathbb{E}[T_4] = \operatorname{tr}\!\big( \symbf{P}\,\ddot{\symbf{\Sigma}}_{kj}\,\symbf{P} \, \symbf{\Sigma} \big).
\end{equation} Using the cyclic property followed by
\(\symbf{P}\symbf{\Sigma}\symbf{P} = \symbf{P}\), \begin{align}
\mathbb{E}[T_4] &= \operatorname{tr}\!\big( \ddot{\symbf{\Sigma}}_{kj}\,\symbf{P} \symbf{\Sigma} \symbf{P} \big) \\
&= \operatorname{tr}\!\big( \ddot{\symbf{\Sigma}}_{kj}\,\symbf{P} \big) \\
&= \operatorname{tr}\!\big( \symbf{P}\,\ddot{\symbf{\Sigma}}_{kj} \big).
\end{align}

\paragraph{Combining the expectations}\label{combining-the-expectations}

Assembling the four components, \begin{align}
\mathbb{E}\!\left[ -\frac{\partial^2 \ell_R}{\partial \theta_k \partial \theta_j} \right]
&= \mathbb{E}[T_1] - \frac12 \mathbb{E}[T_2] + \frac12 \mathbb{E}[T_3] - \frac12 \mathbb{E}[T_4] \\[4pt]
&= \operatorname{tr}\!\big( \symbf{P}\,\dot{\symbf{\Sigma}}_k\,\symbf{P}\,\dot{\symbf{\Sigma}}_j \big)
- \frac12 \operatorname{tr}\!\big( \symbf{P}\,\dot{\symbf{\Sigma}}_j\,\symbf{P}\,\dot{\symbf{\Sigma}}_k \big) \\
&\quad + \frac12 \operatorname{tr}\!\big( \symbf{P}\,\ddot{\symbf{\Sigma}}_{kj} \big)
- \frac12 \operatorname{tr}\!\big( \symbf{P}\,\ddot{\symbf{\Sigma}}_{kj} \big).
\end{align} The two terms involving the second derivatives
\(\ddot{\symbf{\Sigma}}_{kj}\) cancel identically. Moreover, the trace
is symmetric in its matrix arguments, so \begin{equation}
\operatorname{tr}\!\big( \symbf{P}\,\dot{\symbf{\Sigma}}_j\,\symbf{P}\,\dot{\symbf{\Sigma}}_k \big)
= \operatorname{tr}\!\big( \symbf{P}\,\dot{\symbf{\Sigma}}_k\,\symbf{P}\,\dot{\symbf{\Sigma}}_j \big).
\end{equation} Substituting this into the expression above,
\begin{align}
\mathbb{E}\!\left[ -\frac{\partial^2 \ell_R}{\partial \theta_k \partial \theta_j} \right]
&= \operatorname{tr}\!\big( \symbf{P}\,\dot{\symbf{\Sigma}}_k\,\symbf{P}\,\dot{\symbf{\Sigma}}_j \big)
- \frac12 \operatorname{tr}\!\big( \symbf{P}\,\dot{\symbf{\Sigma}}_k\,\symbf{P}\,\dot{\symbf{\Sigma}}_j \big) \\[4pt]
&= \frac12 \operatorname{tr}\!\big( \symbf{P}\,\dot{\symbf{\Sigma}}_k\,\symbf{P}\,\dot{\symbf{\Sigma}}_j \big).
\end{align}

This final expression defines the \((k,j)\)-th entry of the REML Fisher
information matrix, \begin{equation}
\mathcal{I}_{kj}(\symbf{\theta}) = \frac12 \operatorname{tr}\!\big( \symbf{P}\,\dot{\symbf{\Sigma}}_k\,\symbf{P}\,\dot{\symbf{\Sigma}}_j \big).
\end{equation}

\subsubsection{Definition of the Average
Information}\label{definition-of-the-average-information}

The expected information (Fisher information) is \begin{equation}
\mathcal{I}_{kj}(\symbf{\theta}) = \frac12 \operatorname{tr}\!\big( \symbf{P}\,\dot{\symbf{\Sigma}}_k\,\symbf{P}\,\dot{\symbf{\Sigma}}_j \big).
\end{equation} The observed information is denoted
\(\mathcal{J}_{kj}(\symbf{\theta})\), and the \emph{average information}
matrix is defined as \begin{equation}
\mathcal{A}_{kj}(\symbf{\theta}) = \frac12 \Big( \mathcal{I}_{kj}(\symbf{\theta}) + \mathcal{J}_{kj}(\symbf{\theta}) \Big).
\end{equation}

Substituting the expressions and simplifying, the trace terms cancel,
yielding \begin{align}
\mathcal{A}_{kj} &= \frac12 \symbf{Y}^\top \symbf{P}\,\dot{\symbf{\Sigma}}_k\,\symbf{P}\,\dot{\symbf{\Sigma}}_j\,\symbf{P} \symbf{Y} \\
&\quad + \frac14 \operatorname{tr}\!\big( \symbf{P}\,\ddot{\symbf{\Sigma}}_{kj} \big)
- \frac14 \symbf{Y}^\top \symbf{P}\,\ddot{\symbf{\Sigma}}_{kj}\,\symbf{P}\symbf{Y}.
\end{align}

\paragraph{Simplification for Linear Variance Component
Models}\label{simplification-for-linear-variance-component-models}

In the common case where
\(\symbf{\Sigma}(\symbf{\theta}) = \sum_{t=1}^p \theta_t \symbf{V}_t\)
with known matrices \(\symbf{V}_t\), the second derivatives vanish:
\(\ddot{\symbf{\Sigma}}_{kj} = \symbf{0}\) for all \(k,j\). The average
information matrix then reduces to \begin{equation}
\label{eq-AI}
\boxed{
\mathcal{A}_{kj} = \frac12 \,\symbf{Y}^\top \symbf{P}\,\dot{\symbf{\Sigma}}_k\,\symbf{P}\,\dot{\symbf{\Sigma}}_j\,\symbf{P} \symbf{Y}
}.
\end{equation}

This is the form commonly used in the AI-REML algorithm (Gilmour,
Thompson, and Cullis 1995), where the average information matrix is
computed directly from the data without requiring second derivatives.
The algorithm proceeds by iteratively solving

\begin{equation}
\symbf{\theta}^{(t+1)} = \symbf{\theta}^{(t)} + \big( \mathcal{A}(\symbf{\theta}^{(t)}) \big)^{-1} \symbf{S}(\symbf{\theta}^{(t)}).
\end{equation}

\section{Computation of REML}\label{computation-of-reml}

To update \(\theta\), one needs to compute the score function and the AI
matrix, note that the log likelihood function is not necessary. As
derived above, the score function is given in Equation \eqref{eq-score},
and the AI matrix is given in Equation \eqref{eq-AI}. Both have some
common terms that can be shared during computation.

During computation, we will avoid forming \(\symbf{P}\) as it is a
complex matrix, so we will compute combined terms.

\begin{enumerate}
\def\labelenumi{\arabic{enumi}.}
\item
  In both the score and AI matrix, the term \(\symbf{P}\symbf{Y}\)
  commonly appears. Expanding it gives: \begin{equation}
  \symbf{P}\symbf{Y} = \symbf{\Sigma}^{-1}\symbf{Y} - \symbf{\Sigma}^{-1}\symbf{X}(\symbf{X}^\top\symbf{\Sigma}^{-1}\symbf{X})^{-1}\symbf{X}^\top\symbf{\Sigma}^{-1}\symbf{Y}
  \label{eq:PY_expanded}
  \end{equation} Since \(\symbf{P}\) is a residual operator, this is
  actually the weighted residuals. To compute Equation
  \eqref{eq:PY_expanded}, we need to compute: \begin{align}
  &\text{a. } \dot{\symbf{\Sigma}}_k \quad \text{(derivative of covariance matrix w.r.t. $\theta_k$)} \nonumber \\
  &\text{b. } \symbf{L}_{\symbf{\Sigma}} \quad \text{(Cholesky factor)} \nonumber \\
  &\text{c. } \symbf{\Sigma}^{-1} \symbf{Y} \nonumber \\
  &\text{d. } \symbf{\Sigma}^{-1} \symbf{X} \nonumber \\
  &\text{e. } \symbf{X}^\top \symbf{\Sigma}^{-1} \symbf{X} \nonumber \\
  &\text{f. } \symbf{L}_{\symbf{X}^\top \symbf{\Sigma}^{-1} \symbf{X}} \quad \text{(Cholesky factor)} \nonumber \\
  &\text{g. } (\symbf{X}^\top \symbf{\Sigma}^{-1} \symbf{X})^{-1} \symbf{X}^\top \symbf{\Sigma}^{-1} \nonumber \\
  &\text{h. } \symbf{\Sigma}^{-1} \symbf{X} (\symbf{X}^\top \symbf{\Sigma}^{-1} \symbf{X})^{-1} \symbf{X}^\top \symbf{\Sigma}^{-1} \nonumber \\
  &\text{i. } \symbf{\Sigma}^{-1} \symbf{X} (\symbf{X}^\top \symbf{\Sigma}^{-1} \symbf{X})^{-1} \symbf{X}^\top \symbf{\Sigma}^{-1} \symbf{Y} \nonumber
  \end{align} Note that we use Cholesky operations to avoid inverting
  large symmetric matrices.
\item
  Another term that commonly appears is
  \(\symbf{P} \dot{\symbf{\Sigma}}_k\). Expanding it gives:
  \begin{equation}
  \symbf{P} \dot{\symbf{\Sigma}}_k = \symbf{\Sigma}^{-1}\dot{\symbf{\Sigma}}_k - \symbf{\Sigma}^{-1}\symbf{X}(\symbf{X}^\top\symbf{\Sigma}^{-1}\symbf{X})^{-1}\symbf{X}^\top\symbf{\Sigma}^{-1}\dot{\symbf{\Sigma}}_k
  \label{eq:P_Sigma_expanded}
  \end{equation} To compute Equation \eqref{eq:P_Sigma_expanded}, we
  need to compute: \begin{align}
  &\text{a. } \symbf{\Sigma}^{-1} \dot{\symbf{\Sigma}}_k \nonumber \\
  &\text{b. } \symbf{\Sigma}^{-1} \symbf{X} (\symbf{X}^\top \symbf{\Sigma}^{-1} \symbf{X})^{-1} \symbf{X}^\top \symbf{\Sigma}^{-1} \dot{\symbf{\Sigma}}_k \nonumber
  \end{align}
\item
  Then for the score function, we have enough ingredients to compute:
  \begin{equation}
  S_k = \frac{1}{2}\left[\symbf{Y}^\top (\symbf{P} \dot{\symbf{\Sigma}}_k) (\symbf{P}\symbf{Y}) - \operatorname{tr}(\symbf{P} \dot{\symbf{\Sigma}}_k)\right]
  \label{eq:score_computed}
  \end{equation}
\item
  For the AI matrix, we also have enough ingredients to compute it.
  Since the AI matrix is symmetric, we need to compute only half of the
  matrix and then populate it to the full matrix: \begin{equation}
  \mathcal{A}_{kj} = \frac{1}{2}\left[\symbf{Y}^\top (\symbf{P} \dot{\symbf{\Sigma}}_k) (\symbf{P} \dot{\symbf{\Sigma}}_j) (\symbf{P}\symbf{Y})\right]
  \label{eq:AI_computed}
  \end{equation}
\end{enumerate}

During the update, if \(\mathcal{A}\) is large, we can compute
\(\symbf{L}_{\mathcal{A}}\) (the Cholesky factor of \(\mathcal{A}\)) and
then solve \(\mathcal{A}^{-1} \symbf{S}\) efficiently via
back-substitution. Additionally, a damping factor can be applied to
\(\symbf{D} \mathcal{A}^{-1} \symbf{S}\) to control the learning rate.
This is useful when the score or AI matrix computation is inaccurate, or
when the estimation encounters ill-conditioned problems.

\clearpage

\section*{References}\label{references}
\addcontentsline{toc}{section}{References}

\phantomsection\label{refs}
\begin{CSLReferences}{1}{0}
\bibitem[\citeproctext]{ref-gilmour1995average}
Gilmour, Arthur R, Robin Thompson, and Brian R Cullis. 1995. {``Average
Information REML: An Efficient Algorithm for Variance Parameter
Estimation in Linear Mixed Models.''} \emph{Biometrics}, 1440--50.

\end{CSLReferences}




\end{document}
